
\documentclass[11pt,a4paper]{article}

\usepackage[margin=1in]{geometry}
\usepackage[T1]{fontenc}
\usepackage[utf8]{inputenc}
\usepackage{lmodern}
\usepackage{microtype}
\usepackage{amsmath,amssymb}
\usepackage{booktabs}
\usepackage{enumitem}
\usepackage{hyperref}
\usepackage{xcolor}
\usepackage{listings}
\usepackage{fancyhdr}
\usepackage{titlesec}
\usepackage{graphicx}

\hypersetup{
    colorlinks=true,
    linkcolor=blue,
    urlcolor=blue,
    citecolor=blue
}

\definecolor{codegray}{gray}{0.95}

\lstset{
    basicstyle=\ttfamily\small,
    backgroundcolor=\color{codegray},
    frame=single,
    breaklines=true,
    columns=fullflexible,
    keepspaces=true,
    showstringspaces=false
}

\pagestyle{fancy}
\fancyhf{}
\lhead{FRC Research Notes}
\rhead{Week 7}
\cfoot{\thepage}

\titleformat{\section}{\large\bfseries}{\thesection.}{0.5em}{}
\titleformat{\subsection}{\normalsize\bfseries}{\thesubsection.}{0.5em}{}

\title{\textbf{FRC Research --- Week 7 Working Notes}\\
\large GS/MRR Legacy Reproduction and First Hybrid-PIC Tilt Smoke Test}
\author{Juan Pablo Solís Ruiz\\
Supervisor: Prof. Ren Haijun}
\date{October 19--25, 2026}

\begin{document}

\maketitle
\tableofcontents
\newpage

% ============================================================
\section{Purpose and Current Status}

These notes begin a new weekly research record after the completion of the
formation-baseline work and the transition toward the inherited
Grad--Shafranov / Modified Rigid Rotor (GS/MRR) equilibrium and tilt-instability
workflow.

The active objective is now

\[
\boxed{
\text{Reproduce inherited GS/MRR initialization}
\rightarrow
\text{verify provenance}
\rightarrow
\text{run WarpX smoke test}
\rightarrow
\text{measure tilt growth}
}
\]

rather than continuing formation studies.

Progress during the previous two weeks was limited by a high academic workload
and a separate plasma-physics project. No major production simulation was
completed during that interval. The present week therefore resumes the FRC
takeover with a deliberately narrow goal: establish a trustworthy inherited
GS/MRR reference case and verify that it can be evolved with the existing
WarpX build on SCNet.

The guiding principle remains

\[
\boxed{
\text{Reproduce} \rightarrow \text{Verify} \rightarrow \text{Modify}
}
\]

with no physics changes introduced before the inherited case is understood.

% ============================================================
\section{Inherited GS/MRR Materials}

The inherited initialization material currently used for the takeover consists
primarily of:

\begin{lstlisting}
gs_init/
|-- hybrid_init.py
|-- input
`-- modified_rigid_rotorpdf.pdf
\end{lstlisting}

The active script \texttt{hybrid\_init.py} performs the following pipeline:

\[
(\beta_s,\delta_s,n,B_{\rm ext},R_s,\ldots)
\rightarrow
\text{MRR-1 radial closure}
\rightarrow
\text{fixed-boundary GS solve}
\]

\[
\rightarrow
\left\{
\psi(r,z), B_r, B_z, P, n, J_\theta
\right\}
\rightarrow
\text{3D Cartesian interpolation}
\]

\[
\rightarrow
\text{particle loading + } n=1 \text{ tilt seed}
\rightarrow
\texttt{init.bp}
\rightarrow
\text{WarpX input}.
\]

The inherited WarpX input is therefore not an independent hand-written case:
it is generated from the GS/MRR initializer and should be treated together
with the corresponding \texttt{init.bp} state.

% ============================================================
\section{MRR-1 / GS Model Reconstructed}

\subsection{MRR-1 Pressure Closure}

The pressure model implemented in the inherited initializer is

\[
\bar p(x)
=
\beta_s e^{-\alpha x}(1+\sigma x)^n,
\]

with

\[
x =
\frac{\psi}{B_{\rm external}R_s^2}.
\]

For the inherited reference family,

\[
\beta_s = 0.6,
\qquad
n = 2,
\]

and

\[
\delta_s
=
0.4\,\delta_s^{RR}.
\]

The reproduced MRR parameters are

\[
\delta_s = 0.10604579,
\]

\[
\alpha = 31.535058,
\qquad
\sigma = 8.3125473,
\]

\[
\psi_m = -0.053673671.
\]

These values agree with the documented MRR example and are independent of the
absolute magnetic-field scale used in the two cases studied below.

\subsection{Fixed-Boundary Grad--Shafranov Equation}

The initializer solves the axisymmetric FRC equilibrium

\[
\Delta^\star\psi
=
-\mu_0 r^2\frac{dP}{d\psi},
\]

with field reconstruction

\[
B_z
=
\frac{1}{r}\frac{\partial\psi}{\partial r},
\qquad
B_r
=
-\frac{1}{r}\frac{\partial\psi}{\partial z}.
\]

The external axial mirror-field envelope is imposed as a fixed Dirichlet
boundary condition for the GS solve.

The current default initializer uses

\[
B_{\rm external}=0.26~\mathrm{T},
\]

whereas the inherited generated \texttt{input} was found to correspond to

\[
B_{\rm external}=0.13~\mathrm{T}.
\]

This distinction became the main result of the first reproduction audit.

% ============================================================
\section{SCNet Python / openPMD Environment}

The GS/MRR initializer requires \texttt{openpmd\_api} with ADIOS2/BP5 support
in addition to NumPy, SciPy, and Matplotlib.

The working SCNet Python environment is

\begin{table}[htbp]
    \centering
    \caption{Python environment used for the initializer reproduction.}
    \begin{tabular}{ll}
        \toprule
        Component & Version \\
        \midrule
        Python & 3.8.10 \\
        NumPy & 1.24.4 \\
        SciPy & 1.10.1 \\
        Matplotlib & 3.7.5 \\
        openpmd-api & 0.16.1 \\
        \bottomrule
    \end{tabular}
\end{table}

A direct BP5 write/read test completed successfully before the inherited
initializer was executed.

\subsection{Minimal NumPy Compatibility Patch}

The inherited script used

\begin{lstlisting}
np.trapezoid(...)
\end{lstlisting}

in two locations. NumPy 1.24.4 does not provide that API. The working copy was
therefore patched to use

\begin{lstlisting}
np.trapz(...)
\end{lstlisting}

in the two trapezoidal integrations.

No physical or numerical model parameter was changed.

The frozen inherited initializer has SHA256

\begin{lstlisting}
6c8197c175cc7fb4aca6b89f6fde93d7b94bcbac324dd5a161d59b4dadee4034
\end{lstlisting}

and the minimally patched working copy has SHA256

\begin{lstlisting}
4e89d1664cc5ef47bdc610d38eee4bcdcc07e81af5cd72abc25e18aab97a02ae
\end{lstlisting}

The compatibility difference is preserved as
\texttt{numpy\_compat.patch}.

% ============================================================
\section{First Initializer Reproduction: Current 0.26 T Default}

The first successful GS/MRR initializer run used the current default

\[
B_{\rm mid}=0.26~\mathrm{T}.
\]

SCNet job:

\begin{lstlisting}
Job ID: 124911799
State: COMPLETED
Elapsed: 00:01:15
ExitCode: 0:0
Node: a01r1n13
\end{lstlisting}

The initializer reported

\[
I_{\rm target}=391136.95~\mathrm{A},
\]

\[
I_{\rm mid}=391933.99~\mathrm{A},
\]

with

\[
408
\]

nonlinear iterations and final change

\[
4.999\times10^{-5}.
\]

The tilt-seed normalization was

\[
\epsilon=0.01,
\qquad
V_A=134595.59~\mathrm{m/s},
\]

\[
\delta v_{z,\max}=1345.9559~\mathrm{m/s},
\qquad
u_0=4.48651972\times10^{-6}.
\]

The generated initialization dataset was approximately

\[
1.4~\mathrm{GB},
\]

with the principal BP5 payload

\[
\texttt{outputs/init.bp/data.0}
\simeq
1.46\times10^9~\mathrm{bytes}.
\]

% ============================================================
\section{Important Audit Result: 0.26 T Did Not Match the Inherited Input}

The generated 0.26 T WarpX input did not reproduce the inherited input.

The important differences were highly structured rather than arbitrary.

For the inherited input,

\[
n_{0,\rm legacy}
=
4.28012348\times10^{20}~\mathrm{m^{-3}},
\]

whereas the current 0.26 T initializer generated

\[
n_{0,0.26}
=
1.71204939\times10^{21}~\mathrm{m^{-3}}.
\]

Therefore,

\[
\frac{n_{0,0.26}}{n_{0,\rm legacy}}
=
4.
\]

At the same time, the tilt flux normalization changed from approximately

\[
\psi_{O,\rm legacy}
=
3.634298\times10^{-4}
\]

to

\[
\psi_{O,0.26}
=
7.268596\times10^{-4},
\]

a factor of two.

This is consistent with the expected field scaling:

\[
P \propto B^2,
\qquad
n \propto B^2,
\qquad
\psi\propto B.
\]

The Alfv\'en speed remained effectively unchanged,

\[
V_A
=
\frac{B}{\sqrt{\mu_0m_i n}}
\simeq
134595.59~\mathrm{m/s},
\]

because the density scales approximately as \(B^2\).

This immediately suggested that the inherited input was generated using

\[
\boxed{
B_{\rm external}=0.13~\mathrm{T}
}
\]

even though the present initializer default has since been changed to
\(0.26~\mathrm{T}\).

% ============================================================
\section{Legacy 0.13 T Initializer Reproduction}

A second initializer run was therefore launched with the explicit override

\begin{lstlisting}
python hybrid_init.py --b-external 0.13
\end{lstlisting}

SCNet job:

\begin{lstlisting}
Job ID: 124912662
State: COMPLETED
Elapsed: 00:00:51
ExitCode: 0:0
Node: a01r1n05
\end{lstlisting}

The reproduced equilibrium reported

\[
B_{\rm mid}=0.13~\mathrm{T},
\]

\[
T_e=T_i=100~\mathrm{eV},
\qquad
\beta_s=0.6,
\]

\[
I_{\rm target}=195568.47~\mathrm{A},
\]

\[
I_{\rm mid}=195966.99~\mathrm{A}.
\]

The relative current mismatch is approximately

\[
\frac{|I_{\rm mid}-I_{\rm target}|}{I_{\rm target}}
\simeq
2.04\times10^{-3},
\]

or approximately

\[
0.204\%.
\]

The GS solve again required

\[
408
\]

iterations with final convergence measure

\[
4.999\times10^{-5}.
\]

The tilt quantities remained

\[
\epsilon=0.01,
\]

\[
V_A=134595.59~\mathrm{m/s},
\]

\[
\delta v_{z,\max}=1345.9559~\mathrm{m/s},
\]

\[
u_0\simeq4.48652\times10^{-6}.
\]

% ============================================================
\section{Input-Level Reproduction Result}

The generated 0.13 T WarpX input was compared against the frozen inherited
input after ignoring Windows CRLF versus Linux LF line endings.

The only remaining differences were floating-point roundoff in the derived
tilt constants.

For example,

\[
R_0:
1.7067470103522647\times10^{-1}
\rightarrow
1.7067470103588542\times10^{-1},
\]

\[
Z_0:
3.2443343476544395\times10^{-1}
\rightarrow
3.2443343476968189\times10^{-1},
\]

and

\[
\psi_O:
3.6342981445693686\times10^{-4}
\rightarrow
3.6342981445940635\times10^{-4}.
\]

The relative differences are of order

\[
10^{-11}
\]

or smaller.

Therefore the present result is

\[
\boxed{
\text{Legacy 0.13 T seeded MRR/GS initialization reproduced to numerical precision.}
}
\]

This closes the first major uncertainty in the inherited GS/MRR takeover:
the frozen legacy input corresponds to the \(0.13~\mathrm{T}\) case, whereas
the current initializer default corresponds to \(0.26~\mathrm{T}\).

% ============================================================
\section{File-Placement Error and Workflow Correction}

During the first 0.13 T rerun, the copied Slurm script still contained a
hard-coded \texttt{cd} path pointing to the 0.26 T run directory.

As a result, the 0.13 T outputs temporarily overwrote the 0.26 T output
directory.

The simulation itself completed correctly; the issue was only output
placement.

The Slurm script was corrected to use

\begin{lstlisting}
cd "$SLURM_SUBMIT_DIR"
\end{lstlisting}

so future runs write into the directory from which they are submitted.

The reproduced 0.13 T outputs were then moved into the correct run directory.

This reinforces the framework requirement that run directories and provenance
must be generated automatically rather than encoded manually in job scripts.

% ============================================================
\section{First WarpX GS/MRR Tilt Smoke Test}

After reproducing the inherited initializer, a separate smoke-test run
directory was created.

The smoke test uses the reproduced \(0.13~\mathrm{T}\), seeded
\texttt{init.bp} state and the same WarpX executable previously used for the
validated formation work:

\begin{lstlisting}
/public/home/acszlf5t66/WQ/WarpX/build/bin/warpx.3d
\end{lstlisting}

which resolves to

\begin{lstlisting}
warpx.3d.MPI.OMP.DP.PDP.OPMD.EB.QED
\end{lstlisting}

The runtime environment follows the known working formation configuration:

\begin{itemize}
    \item partition: \texttt{kshcnormal},
    \item 1 node,
    \item 4 MPI ranks,
    \item 8 OpenMP threads per rank,
    \item GCC 12.2,
    \item Intel MPI 2021.3.
\end{itemize}

The inherited input contains

\[
N_t=40000,
\qquad
\Delta t=1~\mathrm{ns},
\]

but the smoke test overrides only

\[
\boxed{
N_t=200
}
\]

from the command line.

Therefore

\[
t_{\rm end}
=
200\times1~\mathrm{ns}
=
0.2~\mu\mathrm{s}.
\]

No equilibrium, particle-loading, resistivity, hybrid-substep, or tilt-seed
parameter is modified.

The smoke-test job submitted on October 19 is

\begin{lstlisting}
Job ID: 124913331
Name: frc-tilt-smoke
\end{lstlisting}

The purpose of this run is not yet to measure tilt growth. It tests the full
execution chain

\[
\boxed{
\texttt{init.bp}
\rightarrow
\text{WarpX field/particle loading}
\rightarrow
\text{hybrid-PIC initialization}
\rightarrow
200\text{-step evolution}.
}
\]

At the time these notes were written, the job had been submitted and its final
status was pending.

% ============================================================
\section{Interpretation of the Tilt Seed}

The inherited initializer imposes a controlled axial \(n=1\) perturbation.

The implemented form can be summarized as

\[
\delta u_z
=
u_0
N
\frac{r}{R_0}
\left(
\frac{-\psi}{\psi_O}
\right)^p
\cos(\theta-\phi),
\qquad
\psi<0,
\]

with

\[
p=2,
\qquad
\phi=0,
\qquad
\epsilon=0.01.
\]

The perturbation is zero outside the closed-flux region selected by
\(\psi<0\).

Importantly, this seed is already written into the particle momenta stored in
\texttt{init.bp}. The constants written into the WarpX input are therefore
provenance information; changing them without regenerating
\texttt{init.bp} would not change the actual initialized particle
perturbation.

% ============================================================
\section{What Has Been Established So Far}

The following statements are now supported by the reproduction work:

\begin{enumerate}
    \item The inherited MRR-1 / fixed-boundary GS initializer runs successfully
    on SCNet after a minimal NumPy API compatibility patch.

    \item The current initializer default is a \(0.26~\mathrm{T}\) case.

    \item The frozen inherited WarpX input corresponds instead to the
    \(0.13~\mathrm{T}\) case.

    \item The \(0.13~\mathrm{T}\) generated input reproduces the inherited input
    to numerical precision after normalizing line endings.

    \item The expected scalings
    \[
    n\propto B^2,
    \qquad
    \psi\propto B
    \]
    explain the difference between the 0.13 T and 0.26 T generated cases.

    \item The Alfv\'en speed and normalized tilt-seed velocity remain
    essentially unchanged between these two scaled equilibria.

    \item A first WarpX hybrid-PIC smoke test has been submitted from the
    reproduced \(0.13~\mathrm{T}\) state.
\end{enumerate}

% ============================================================
\section{Immediate Next Steps}

The next actions should be performed in the following order.

\begin{enumerate}
    \item Check the final state of job \texttt{124913331}.

    \item If the 200-step smoke test completes successfully, verify:
    \begin{itemize}
        \item WarpX used-input provenance,
        \item successful field and external-particle loading,
        \item diagnostic output at the expected intervals,
        \item absence of NaNs or immediate numerical failure.
    \end{itemize}

    \item Build a first quantitative \(n=1\) tilt diagnostic.

    A useful geometric metric is the complex axial displacement moment

    \[
    Z_1(t)
    =
    \frac{
        \int z\,\rho(\mathbf{x},t)e^{-i\theta}\,dV
    }{
        \int \rho(\mathbf{x},t)\,dV
    },
    \]

    with

    \[
    A_1(t)=|Z_1(t)|.
    \]

    \item Once the smoke test is trusted, run a longer seeded case long enough
    to identify a linear-growth interval.

    \item Fit

    \[
    A_1(t)\propto e^{\gamma t}
    \]

    or equivalently

    \[
    \ln A_1(t)=\gamma t+C
    \]

    to obtain the dimensional tilt growth rate.

    \item Only after reproducing the inherited seeded case should the study move
    to seed-amplitude scans, unseeded cases, \(B_{\rm ext}\) scaling, or other
    equilibrium modifications.
\end{enumerate}

% ============================================================
\section{Open Questions for the Next Audit}

Several normalization and diagnostic questions remain intentionally open.

\begin{itemize}
    \item What exact length scale was used in the inherited definition of
    \(\tau_A\)?

    \item Which evolved geometry definition was used for \(R_s\), \(L_s\),
    elongation \(E\), and \(S^*=R_s/d_i\)?

    \item Should the primary tilt amplitude be based on density displacement,
    magnetic-axis displacement, or a direct azimuthal Fourier decomposition?

    \item Does the inherited post-processing use ion bulk velocity, total
    current, or an approximation such as \(j_z/\rho\) when diagnosing the tilt?

    \item Which inherited run should be treated as the canonical reference:
    the reproduced \(0.13~\mathrm{T}\) case or one of the later
    \(0.26~\mathrm{T}\) seeded/unseeded cases?
\end{itemize}

These questions should be resolved through reproduction and explicit
diagnostics rather than assumptions.

% ============================================================
\section{End-of-Day Status --- October 19}

The main result of the day is not a new parameter scan. It is the establishment
of a trustworthy GS/MRR reference initialization.

The sequence completed was

\[
\boxed{
\text{legacy audit}
\rightarrow
\text{SCNet dependency setup}
\rightarrow
\text{0.26 T reproduction}
\rightarrow
\text{field-scaling discrepancy}
\rightarrow
\text{0.13 T reproduction}
}
\]

followed by

\[
\boxed{
\text{first WarpX GS/MRR smoke test submitted}.
}
\]

This is the first end-to-end step from the inherited equilibrium generator
toward a reproducible tilt-instability workflow.

% ============================================================
\section{October 19, 2026 --- Legacy GS/MRR Reproduction and
ADIOS2-Enabled WarpX Takeover}
% ============================================================

Today the main objective was to move from inspection of the inherited
Grad--Shafranov / Modified Rigid Rotor (GS/MRR) initialization toward an
actual reproducible tilt-instability simulation. The work resulted in two
important milestones: (i) numerical reproduction of the inherited
$B_{\mathrm{ext}}=0.13~\mathrm{T}$ equilibrium initialization, and
(ii) a working local MPI-enabled ADIOS2 installation suitable for rebuilding
WarpX with BP5 support.

% ------------------------------------------------------------
\subsection{Legacy GS/MRR initialization reproduced}
% ------------------------------------------------------------

The inherited hybrid initializer was first reproduced with the current script
default,
\[
B_{\mathrm{ext}} = 0.26~\mathrm{T},
\]
using only a minimal NumPy compatibility modification
(\texttt{np.trapezoid} $\rightarrow$ \texttt{np.trapz}). No physical or
numerical parameters were changed.

The $0.26~\mathrm{T}$ run converged after 408 outer iterations with
\[
I_{\mathrm{target}} \simeq 391.14~\mathrm{kA},
\qquad
I_{\mathrm{mid}} \simeq 391.93~\mathrm{kA},
\]
and a final GS relative change of approximately
\[
5.0\times10^{-5}.
\]

Comparison against the inherited WarpX input revealed a systematic scaling:
the generated density was four times larger while the characteristic
poloidal flux was twice as large,
\[
n \propto B^2,
\qquad
\psi \propto B.
\]
At the same time, the Alfv\'en speed remained essentially unchanged because
\[
V_A \sim \frac{B}{\sqrt{\mu_0 m_i n}}
\]
is invariant under the simultaneous scaling
$B\rightarrow2B$ and $n\rightarrow4n$.

This strongly indicated that the inherited input had originally been
generated with
\[
B_{\mathrm{ext}} = 0.13~\mathrm{T},
\]
consistent with the MRR documentation.

The initializer was therefore rerun at $0.13~\mathrm{T}$. The resulting
equilibrium gave
\[
I_{\mathrm{target}} \simeq 195.57~\mathrm{kA},
\qquad
I_{\mathrm{mid}} \simeq 195.97~\mathrm{kA},
\]
again converging in 408 iterations.

The regenerated WarpX input matched the inherited input to numerical
precision. The only differences were floating-point roundoff in quantities
such as
\[
u_0,\quad R_0,\quad Z_0,\quad \psi_0,\quad \text{and}\quad N_{\mathrm{tilt}},
\]
with relative discrepancies of order
\[
10^{-11}
\]
or smaller.

Therefore, the main equilibrium-reproduction result for today is

\begin{quote}
\textbf{The inherited $0.13~\mathrm{T}$ seeded MRR/GS initialization was
reproduced to numerical precision.}
\end{quote}

The regenerated equilibrium also reproduces the inherited $n=1$ tilt seed.
For the current initialization,
\[
\epsilon = 0.01,
\qquad
V_A \simeq 1.346\times10^5~\mathrm{m\,s^{-1}},
\]
giving a peak imposed axial perturbation of approximately
\[
\delta u_z^{\mathrm{max}}
\simeq 1.346\times10^3~\mathrm{m\,s^{-1}}.
\]

% ------------------------------------------------------------
\subsection{First inherited tilt smoke test}
% ------------------------------------------------------------

A short 200-step WarpX smoke test was prepared using the reproduced
$0.13~\mathrm{T}$ initial condition. The test used the same inherited input
physics while reducing only the total number of steps through the command-line
override
\[
N_{\mathrm{step}} = 200.
\]

The existing WarpX executable initialized MPI and OpenMP successfully, but
terminated before advancing the physical simulation. The failure occurred
when openPMD attempted to read the BP5 initialization file:

\begin{verbatim}
openPMD::error::WrongAPIUsage
openPMD-api built without support for backend 'ADIOS2'
\end{verbatim}

Thus, this failure was identified as an I/O/build-infrastructure problem
rather than a failure of the GS equilibrium or tilt physics.

Inspection of the original WarpX CMake configuration confirmed that
\texttt{ADIOS2\_DIR} was not found when the existing WarpX/openPMD executable
was compiled.

% ------------------------------------------------------------
\subsection{Local ADIOS2 build for BP5 support}
% ------------------------------------------------------------

A separate ADIOS2 installation was built inside the FRC sandbox, leaving the
existing WarpX installation untouched.

The final build used ADIOS2 v2.11.0 with source commit
\begin{verbatim}
a3eb1ab7d713165f457377a925c0c4f3f4dcf0e5
\end{verbatim}

and was configured approximately as
\begin{verbatim}
MPI          = ON
SST          = OFF
HDF5         = OFF
OpenSSL      = OFF
Python       = OFF
Fortran      = OFF
Shared libs  = OFF
\end{verbatim}

Two cluster-specific build issues were isolated during this process.

First, the bundled FFS sources were resolving an incorrect third-party
\texttt{config.h}. The generated FFS configuration itself was correct, but
source files were instead finding the DILL configuration header. A local
build-only workaround was used to ensure that the FFS sources resolved their
own generated configuration. This allowed the FFS target to build
successfully.

Second, EVPath was still being built even though SST had explicitly been
disabled. Since the present BP5 workflow does not require SST, a minimal
CMake patch was introduced so that EVPath is only configured when
\texttt{ADIOS2\_HAVE\_SST} is enabled. The patch was saved separately for
provenance:

\begin{verbatim}
adios2-no-evpath-without-sst.patch
SHA256:
04278ae5f6522401d69909b94da68988b324781393794c701deceae70c669996
\end{verbatim}

With these build-only modifications, ADIOS2 compiled and installed
successfully:

\begin{verbatim}
Slurm job:   124923692
State:       COMPLETED
Exit code:   0
Elapsed:     00:03:45
\end{verbatim}

The completed build explicitly produced the FFS and ADIOS2 core targets,
including the BP5 reader/writer implementation, and installed the expected
static libraries such as

\begin{verbatim}
libadios2_core.a
libadios2_core_mpi.a
libadios2_cxx.a
libadios2_cxx_mpi.a
libadios2_c.a
libadios2_c_mpi.a
\end{verbatim}

This closes the ADIOS2 dependency needed to rebuild openPMD/WarpX with BP5
support.

% ------------------------------------------------------------
\subsection{New WarpX build and dependent smoke test}
% ------------------------------------------------------------

To avoid modifying the inherited WarpX installation, a clean WarpX source
snapshot and a clean openPMD-api snapshot are being used for a separate
build. The new build is configured to use the locally installed ADIOS2 and
to preserve the important physics/numerics settings of the previous
executable:

\begin{align*}
\text{geometry} &:\quad 3\mathrm{D},\\
\text{parallelism} &:\quad \mathrm{MPI}+\mathrm{OpenMP},\\
\text{field precision} &:\quad \mathrm{double},\\
\text{particle precision} &:\quad \mathrm{double},\\
\text{openPMD} &:\quad \mathrm{ON},\\
\text{ADIOS2} &:\quad \mathrm{ON},\\
\text{EB} &:\quad \mathrm{ON},\\
\text{QED} &:\quad \mathrm{ON}.
\end{align*}

At the end of today's work, the following Slurm dependency chain was active:

\begin{verbatim}
WarpX + openPMD + ADIOS2 build:
    Job 124925297

200-step tilt smoke test:
    Job 124925300

Dependency:
    afterok:124925297
\end{verbatim}

Therefore, the smoke test will execute automatically only if the new WarpX
build completes successfully.

\textbf{Current status:} the $0.13~\mathrm{T}$ GS/MRR initialization is now
reproduced and frozen as the inherited reference case; ADIOS2/BP5 support has
been successfully built; the ADIOS2-enabled WarpX executable and the first
200-step physics smoke test are still pending and should not yet be reported
as successful.

% ------------------------------------------------------------
\subsection{Immediate next step}
% ------------------------------------------------------------

If the dependent smoke test successfully reads \texttt{init.bp}, advances
200 steps, and writes diagnostics, the next stage will be to extend the
simulation into the early tilt-growth regime and construct an $n=1$
displacement diagnostic, for example
\[
Z_1(t)
=
\frac{
\displaystyle
\int z\,\rho(\mathbf{x},t)e^{-i\theta}\,dV
}{
\displaystyle
\int \rho(\mathbf{x},t)\,dV
},
\qquad
A_1(t)=|Z_1(t)|.
\]

During a linear-growth interval,
\[
A_1(t)\propto e^{\gamma t},
\]
so that
\[
\ln A_1(t)=\gamma t+C
\]
can be used to estimate the physical $n=1$ tilt growth rate before proceeding
to seed-amplitude and parameter scans.

\end{document}
